\section{The Dirac field with mass term and self-interaction}
\label{sec:dirac}

\subsection{Matrices and spin connection}

At each point of a curved spacetime the tangent space is flat. The Dirac matrices of the orthonormal
frame are therefore eight constant matrices $\gamma^{A}$, $A=0,\dots,7$. They are real $16\times16$
matrices with entries $0$ and $\pm1$, printed in full in Appendix~\ref{app:matrices}.
%%<1644>
The four matrices of the space-like coordinates are symmetric and square to $+1$, the four matrices
of the time-like coordinates are antisymmetric and square to $-1$, and different matrices
anticommute:
\begin{equation}
\gamma^{A}\gamma^{B}+\gamma^{B}\gamma^{A}=2\,\etamup{AB} .
\label{eq:clifford}
\end{equation}
%%</1644>
%%<1671>
The matrix of the time, $\gamma^{4}$, is symmetric and squares to $+1$; the seven matrices of the
space-like coordinates are antisymmetric and square to $-1$; different matrices anticommute:
\begin{equation}
\gamma^{A}\gamma^{B}+\gamma^{B}\gamma^{A}=-2\,\etamup{AB} .
\label{eq:clifford}
\end{equation}
The minus sign is a consequence of the sign convention of the metric (space-like components
positive) together with matrices whose time matrix squares to $+1$; it accompanies the matrices
through all formulas below.
%%</1671>
All of these properties were verified by exact integer arithmetic for this work. The matrices with a
coordinate index are $\gamma^{\mu}=e_A{}^{\mu}\gamma^{A}$: five of them are the constant matrices,
and the three along the deflating coordinates are the constant matrices divided by the scale
factors, $\gamma^{A=t}/b_t$. They satisfy
$\gamma^{\mu}\gamma^{\nu}+\gamma^{\nu}\gamma^{\mu}=\SIG{2}{-2}\,\gmup{\mu\nu}$ with the inverse
metric $\gmup{\mu\nu}$.

A spinor is defined relative to a frame, and its derivative must account for the turning of the frame
from point to point \cite{BrillWheeler1957}. The covariant derivative of a spinor is
$\nabla_\mu\psi=(\partial_\mu+\Omega_\mu)\psi$ with the spin connection
\begin{equation}
\begin{aligned}
\Omega_\mu&=\SIG{}{-}\tfrac14\,\omega_{\mu AB}\,\gamma^{A}\gamma^{B},\\
\omega_{\mu AB}&=\etam_{AC}\,e^{C}{}_{\nu}\big(\partial_\mu e_B{}^{\nu}+\Gamma^{\nu}{}_{\mu\lambda}e_B{}^{\lambda}\big),
\end{aligned}
\label{eq:spindef}
\end{equation}
whose coefficients $\omega_{\mu AB}=-\omega_{\mu BA}$ follow from the metric ($e^{C}{}_{\nu}$ is
the inverse of the frame). \SIG{}{The minus sign is that of Eq.~(\ref{eq:clifford}): the matrices
$-\tfrac14[\gamma^{A},\gamma^{B}]$ generate the Lorentz transformations of the spinors. }For the
metric (\ref{eq:metric}) the only coefficients that do not vanish are
$\omega_{t\,t4}=-\omega_{t\,4t}=\SIG{}{-}\rt\,b_t$, so that
\begin{equation}
\begin{aligned}
\Omega_t&=\tfrac12\,\rt\,b_t\,\gamma^{t}\gamma^{4}\quad(t=5,6,7),\\
\Omega_0&=\dots=\Omega_4=0,
\end{aligned}
\label{eq:spinconn}
\end{equation}
with the constant matrices of the frame on the right. Each $\Omega_t$ generates the Lorentz
transformation in the plane of $xt$ and $x4$. The spin connection is correct if and only if the
Dirac matrices are covariantly constant,
$\partial_\mu\gamma^{\nu}+\Gamma^{\nu}{}_{\mu\lambda}\gamma^{\lambda}+[\Omega_\mu,\gamma^{\nu}]=0$,
which holds for all 64 pairs of indices. Its curvature is the Riemann tensor,
$[\nabla_\mu,\nabla_\nu]\psi=\SIG{}{-}\tfrac14R_{AB\mu\nu}\gamma^{A}\gamma^{B}\psi$, a third route to
Eq.~(\ref{eq:riemann}). One combination enters the Dirac equation:
\begin{equation}
\gamma^{\mu}\Omega_\mu=-\tfrac12\,\gsum\,\gamma^{4} .
\label{eq:gammaomega}
\end{equation}
The scale factors cancel, and only the total rate remains.

\subsection{The Lagrangian: kinetic term, mass term, self-interaction}
\label{sec:lagrangian}

The field of this paper, $\psi_{\tg}(x)$, is a column of sixteen complex-valued Grassmann operators.
Its Lagrangian has three parts,
\begin{equation}
\mathcal L_{\tg}=\Vd\,\big[\mathcal L_{\mathrm{kin}}+\mathcal L_{\mathrm{mass}}+\Lself\big],
\label{eq:Ltot}
\end{equation}
where $\Vd=\sqrt{|\det g|}$ is the factor that makes $\mathcal L_{\tg}\,d^{8}x$ invariant. A
bosonic field theory would call the second part the ``mass times $\phi^{2}$'' term and the third the
``$\phi^{4}$'' term; for a fermion field they are bilinear and quartic in $\psi$.

%%<1671>
\emph{Kinetic and mass terms.} With the adjoint $\bar\psi=\psi^\dagger\gamma^{4}$,
\begin{equation}
\mathcal L_{\mathrm{kin}}+\mathcal L_{\mathrm{mass}}=\bar\psi\,\big(i\gamma^{\mu}\nabla_\mu-m\big)\,\psi .
\label{eq:Lfree}
\end{equation}
The mass term $-m\,\bar\psi\psi$ is real because $\gamma^{4}$ is symmetric, and it is a scalar under
the Lorentz transformations of the frame.

\emph{Self-interaction.} The most general interaction without derivatives that is quartic in the
field and invariant under the Lorentz transformations of the frame is built from the products of
the Dirac matrices. Let $K=(A_1<A_2<\dots<A_r)$ be a strictly increasing list of $r$ indices out of
$0,\dots,7$; $r=|K|$ is its \emph{rank}, and there are $2^{8}=256$ lists, with
$1,8,28,56,70,56,28,8,1$ of the ranks $0,\dots,8$. Define
\begin{equation}
\Theta^{K}=\zeta_r\,\gamma^{A_1}\gamma^{A_2}\cdots\gamma^{A_r},\qquad
\zeta_r=\begin{cases}i,&r=2,3,6,7,\\1,&\text{otherwise},\end{cases}
\label{eq:Theta}
\end{equation}
($\Theta^{()}=1$ for the empty list) and $\Theta_K=(\Theta^{K})^{-1}=\sgnK\,\Theta^{K}$, where the
sign $\sgnK=\pm1$ is the product of the squares of the matrices in $K$. The phase $\zeta_r$ makes
$\gamma^{4}\Theta^{K}$ Hermitian, so that every bilinear $\bar\psi\Theta^{K}\psi$ is real. The
\emph{invariant of rank $r$} and the self-interaction are
\begin{equation}
I_r=\sum_{|K|=r}(\bar\psi\Theta_K\psi)(\bar\psi\Theta^{K}\psi),\qquad
\Lself=-\frac12\sum_{r=0}^{8}g_r\,I_r .
\label{eq:Lself}
\end{equation}
The nine real constants $g_r$ are the couplings. $I_0=(\bar\psi\psi)^{2}$ is the square of the scalar
density and $I_1=\sum_A(\bar\psi\gamma_{A}\psi)(\bar\psi\gamma^{A}\psi)$ the square of the current, where
$\gamma_A=(\gamma^{A})^{-1}$; with the chirality matrix $\chi$ of Sec.~\ref{sec:clock},
$I_8=-(\bar\psi\,i\chi\,\psi)^{2}$ and
$I_7=-\sum_A(\bar\psi\gamma_{A}\chi\psi)(\bar\psi\gamma^{A}\chi\psi)$. In the traditional names the scalar,
vector, axial-vector and pseudoscalar couplings are $g_S=g_0$, $g_V=g_1$, $g_A=-g_7$, $g_P=-g_8$. A
coupling $g_r$ carries one digit as its subscript; it must not be confused with a rate ($\gff$) or
with a component of the metric ($\gm_{55}$). In a gravitational field each $I_r$ is written with the
frame matrices exactly as in flat spacetime, because it is invariant under the Lorentz
transformations of the frame. Section~\ref{sec:selfint} shows that for anticommuting fields only
four of the nine invariants are independent.

\subsection{The field equation}

The variation of the action with respect to $\bar\psi$ gives
\begin{equation}
\big(i\gamma^{\mu}\nabla_\mu-m\big)\psi_{\tg}=\sum_{r=0}^{8}g_r\sum_{|K|=r}\big(\bar\psi\Theta_K\psi\big)\,\Theta^{K}\psi .
\label{eq:dirac}
\end{equation}
(The bilinears are even elements of the Grassmann algebra, so that the derivative with respect to
$\bar\psi$ acts on $\Lself$ as on a product of ordinary numbers.) The left-hand side is the Dirac
operator with the mass term; the right-hand side is the self-interaction, cubic in the field.
Without it the equation is linear. With Eq.~(\ref{eq:gammaomega}),
\begin{equation}
\begin{aligned}
&i\gamma^{4}\Big(\partial_4-\frac{\gsum}{2}\Big)\psi+i\sum_{j=0}^{3}\gamma^{j}\partial_j\psi
+i\sum_{t=5}^{7}\frac{\gamma^{t}}{b_t}\,\partial_t\psi-m\,\psi\\
&\qquad=\sum_{r=0}^{8}g_r\sum_{|K|=r}\big(\bar\psi\Theta_K\psi\big)\,\Theta^{K}\psi,
\end{aligned}
\label{eq:diracexp}
\end{equation}
%%</1671>
%%<1644>
\emph{Kinetic and mass terms.} With the adjoint $\tilde\psi=\psi^\dagger\Cm$, where
$\Cm=\gamma^{0}\gamma^{1}\gamma^{2}\gamma^{3}$ (a real symmetric matrix that turns each Dirac matrix
into minus its transpose),
\begin{equation}
\mathcal L_{\mathrm{kin}}+\mathcal L_{\mathrm{mass}}=\tilde\psi\,\big(\gamma^{\mu}\nabla_\mu-m\big)\,\psi .
\label{eq:Lfree}
\end{equation}
It contains no factor $i$: with the real matrices of $4+4$ the operator $\gamma^{\mu}\nabla_\mu-m$
gives the mass shell of a particle of mass $m$. The mass term $-m\,\tilde\psi\psi$ is real because
$\Cm$ is symmetric, and it is invariant under the transformations of $SO(4,4)$ of the frame.

\emph{Self-interaction.} The most general interaction without derivatives that is quartic in the
field and invariant under the transformations of the frame is built from the products of the Dirac
matrices. Let $K=(A_1<A_2<\dots<A_r)$ be a strictly increasing list of $r$ indices out of
$0,\dots,7$; $r=|K|$ is its \emph{rank}, and there are $2^{8}=256$ lists, with
$1,8,28,56,70,56,28,8,1$ of the ranks $0,\dots,8$. Define
\begin{equation}
\Theta^{K}=\zeta_r\,\gamma^{A_1}\gamma^{A_2}\cdots\gamma^{A_r},\qquad
\zeta_r=\begin{cases}i,&r=1,2,5,6,\\1,&\text{otherwise},\end{cases}
\label{eq:Theta}
\end{equation}
($\Theta^{()}=1$ for the empty list) and $\Theta_K=(\Theta^{K})^{-1}=\sgnK\,\Theta^{K}$, where the
sign $\sgnK=\pm1$ is $(-1)^{r}$ times the product of the squares of the matrices in $K$. The phase
$\zeta_r$ makes $\Cm\Theta^{K}$ Hermitian, so that every bilinear $\tilde\psi\Theta^{K}\psi$ is real.
The \emph{invariant of rank $r$} and the self-interaction are
\begin{equation}
I_r=\sum_{|K|=r}(\tilde\psi\Theta_K\psi)(\tilde\psi\Theta^{K}\psi),\qquad
\Lself=-\frac12\sum_{r=0}^{8}g_r\,I_r .
\label{eq:Lself}
\end{equation}
The nine real constants $g_r$ are the couplings: $g_0$ multiplies the square of the scalar
$\tilde\psi\psi$, $g_1$ the square of the current $\tilde\psi\gamma^{A}\psi$, $g_7$ and $g_8$ the
squares of the axial current and of the pseudoscalar; in the traditional names $g_S=g_0$, $g_V=g_1$,
$g_A=-g_7$, $g_P=-g_8$. A coupling $g_r$ carries one digit as its subscript; it must not be confused
with a rate ($\gff$) or with a component of the metric ($\gm_{55}$). In a gravitational field each
$I_r$ is written with the frame matrices exactly as in flat spacetime, because it is invariant under
the transformations of the frame. Equation~(\ref{eq:Lself}) is the \emph{covariant form} of the
self-interaction, called \emph{form C} below; it is the definition of the self-interaction of the
$4+4$ model in this paper. Section~\ref{sec:selfint} shows that for anticommuting fields only four of
the nine invariants are independent. Section~\ref{sec:forms} shows that the kinetic term
(\ref{eq:Lfree}) is indefinite in the sector of the theory that can be quantized, defines the model
that is quantized (the positive kinetic term with the self-interaction of form C), and records a
second form of the self-interaction, \emph{form Q}, as the alternative.

\subsection{The field equation}

The variation of the action with respect to $\tilde\psi$ gives
\begin{equation}
\big(\gamma^{\mu}\nabla_\mu-m\big)\psi_{\tg}=\sum_{r=0}^{8}g_r\sum_{|K|=r}\big(\tilde\psi\Theta_K\psi\big)\,\Theta^{K}\psi .
\label{eq:dirac}
\end{equation}
(The bilinears are even elements of the Grassmann algebra, so that the derivative with respect to
$\tilde\psi$ acts on $\Lself$ as on a product of ordinary numbers.) The left-hand side is the Dirac
operator with the mass term; the right-hand side is the self-interaction, cubic in the field.
Without it the equation is linear. With Eq.~(\ref{eq:gammaomega}),
\begin{equation}
\begin{aligned}
&\gamma^{4}\Big(\partial_4-\frac{\gsum}{2}\Big)\psi+\sum_{j=0}^{3}\gamma^{j}\partial_j\psi
+\sum_{t=5}^{7}\frac{\gamma^{t}}{b_t}\,\partial_t\psi-m\,\psi\\
&\qquad=\sum_{r=0}^{8}g_r\sum_{|K|=r}\big(\tilde\psi\Theta_K\psi\big)\,\Theta^{K}\psi,
\end{aligned}
\label{eq:diracexp}
\end{equation}
%%</1644>
with the constant matrices. Three things distinguish the left-hand side and the interaction from
their forms on flat spacetime: the derivative along the time is shifted by $\gsum/2$, each
derivative along a deflating coordinate is divided by its scale factor, and, as the next subsection
shows, the self-interaction is strengthened by the shrinking of the volume.

\subsection{The weighted field}

The shift is removed by a change of normalization. Define the \emph{weighted field}
\begin{equation}
\hpsi=\sqrt{\Vd}\,\psi .
\label{eq:weighted}
\end{equation}
Because $\partial_4\sqrt{\Vd}=-\tfrac12\gsum\sqrt{\Vd}$, the left-hand side of
Eq.~(\ref{eq:diracexp}) applied to $\psi$ equals $\Vd^{-1/2}$ times the same expression without the
term $\gsum/2$ applied to $\hpsi$. The right-hand side is cubic in the field and therefore equals
$\Vd^{-3/2}$ times the same expression in $\hpsi$. The weighted field thus obeys
Eq.~(\ref{eq:diracexp}) \emph{without the term $\gsum/2$ and with every coupling $g_r$ replaced by}
\begin{equation}
\frac{g_r}{\Vd(x4)}\ \propto\ g_r\,e^{\gsum\,x4} .
\label{eq:grun}
\end{equation}
The quantity $\psi^\dagger\psi$ is a density per unit of proper volume and grows as the deflating
coordinates shrink; $\hpsi^\dagger\hpsi$ is the density per unit of coordinate volume, from which
this trivial growth is removed.

The central observation of this paper follows at once. \emph{For a field that does not depend on
$x5$, $x6$, $x7$, the equation of the weighted field is exactly the flat equation with a
self-interaction whose couplings grow as $e^{\gsum\,x4}$.} Every result that holds on flat
spacetime for such fields, in particular everything about the pairing of the fermions, holds for
the weighted field at each moment with the couplings of that moment. The deflation enters those
results in one way only, through the dependence of the couplings on $\Vd$ (Sec.~\ref{sec:clock}):
the mass term is not changed, and every quartic coupling grows at the same rate.

A field that does depend on a deflating coordinate feels the deflation directly. The rest of this
section and Sec.~\ref{sec:modes} treat the equation without the self-interaction, which is linear
and can be solved exactly; it describes a field so dilute that the cubic term is negligible, and
the single-particle states from which the many-body treatment of
Secs.~\ref{sec:selfint}--\ref{sec:ks} starts. Because the
coefficients depend on $x4$ only, a solution of the linear equation can be taken proportional to
$\exp(i\sum_j\SIG{k_j}{p_j}\,xj+i\sum_tk_t\,xt)$ with seven constant wavenumbers, four along the flat
coordinates $xj$, $j=0,1,2,3$, and three along the deflating ones; such a solution is a \emph{mode}. Multiplying the equation of the weighted field by $\gamma^{4}$ gives an ordinary
differential equation in the time:
%%<1644>
\begin{equation}
\partial_4\hpsi=K(x4)\,\hpsi,\qquad K=K_\perp+i\,\kd(x4)\,\gamma^{4}\gamma^{d},
\label{eq:K}
\end{equation}
where $K_\perp=-m\gamma^{4}+i\sum_jk_j\gamma^{4}\gamma^{j}$ is the operator of the modes that do not
depend on the deflating times,
\begin{equation}
\kd(x4)=\Big(\sum_t\frac{k_t^{2}}{b_t(x4)^{2}}\Big)^{1/2}
\label{eq:kd}
\end{equation}
is the \emph{physical} wavenumber along the deflating coordinates, and
$\gamma^{d}=\kd^{-1}\sum_t(k_t/b_t)\gamma^{t}$ is the frame matrix along its direction. With
$\Eperp=(m^{2}+\sum_jk_j^{2})^{1/2}$,
\begin{equation}
K^{2}=\big(\kd^{2}-\Eperp^{2}\big)\,1 .
\label{eq:K2}
\end{equation}
$K_\perp$ is anti-Hermitian and $i\gamma^{4}\gamma^{d}$ is Hermitian. The evolution in the time is
unitary, and $\hpsi^\dagger\hpsi$ is conserved, if and only if $k_5=k_6=k_7=0$. For the other modes
the square of $K$ is negative while $\kd<\Eperp$ and positive afterwards: the mode oscillates until
its physical wavenumber, which grows as $e^{\rt\,x4}$, has reached $\Eperp$, and its solutions grow
or decay exponentially thereafter. A conserved charge exists, $\hpsi^\dagger W_{1644}\hpsi$ with the
Hermitian matrix $W_{1644}=i\Cm\gamma^{4}$ ($W_{1644}^{2}=1$), but it is not positive.
%%</1644>
%%<1671>
\begin{equation}
i\,\partial_4\hpsi=\hh(x4)\,\hpsi,\qquad\hh=\hh_\perp+\kd(x4)\,\alpha^{d},
\label{eq:h}
\end{equation}
where $\hh_\perp=m\beta+\sum_jp_j\alpha^{j}$ is the Hamiltonian of the modes that do not depend on
the deflating coordinates ($\beta=\gamma^{4}$, $\alpha^{A}=\gamma^{4}\gamma^{A}$),
\begin{equation}
\kd(x4)=\Big(\sum_t\frac{k_t^{2}}{b_t(x4)^{2}}\Big)^{1/2}
\label{eq:kd}
\end{equation}
is the \emph{physical} wavenumber along the deflating coordinates, and
$\alpha^{d}=\kd^{-1}\sum_t(k_t/b_t)\alpha^{t}$ is the matrix $\alpha$ along its direction. With
$\Eperp=(m^{2}+\sum_jp_j^{2})^{1/2}$,
\begin{equation}
\hh^{2}=\big(\Eperp^{2}+\kd^{2}\big)\,1 .
\label{eq:h2}
\end{equation}
The Hamiltonian is real and symmetric at every time: the evolution is unitary and
$\hpsi^\dagger\hpsi$ is conserved. But the Hamiltonian depends on the time, through a term that
anticommutes with $\beta$, and the energy of a mode with $k_t\neq0$ grows as $e^{\rt\,x4}$. A
Hamiltonian of this kind creates pairs \cite{Parker1969}.
%%</1671>

